\documentclass{resume} % Use the custom resume.cls style
\usepackage[T1]{fontenc}
\usepackage[utf8]{inputenc}
\usepackage{cmap}
\input{glyphtounicode}
\pdfgentounicode=1
\usepackage[dvipsnames]{xcolor}
\definecolor{LinkNavy}{HTML}{12324A}
\usepackage{hyperref}
\hypersetup{
    colorlinks=true,
    linkcolor=LinkNavy,
    citecolor=black,
    urlcolor=LinkNavy
}
\usepackage{enumitem}
\usepackage[backend=bibtex, style=numeric, sorting=ydnt, maxnames=99, maxcitenames=99]{biblatex}
\addbibresource{ref.bib}
\usepackage{palatino}
\usepackage[left=0.75in,top=0.75in,right=0.75in,bottom=0.75in]{geometry} % Document margins
\usepackage{fancyhdr}
\pagestyle{fancy}
\fancyhf{}
\fancyfoot[C]{\thepage}
\renewcommand{\headrulewidth}{0pt}

\name{Wei Guo}
\address{Personal website: \href{https://alexandreguo2001.github.io/}{\texttt{alexandreguo2001.github.io}}}
\address{
    \href{https://scholar.google.com/citations?user=0r_TzzkAAAAJ&hl=en}{Google Scholar} \\ 
    \href{https://www.semanticscholar.org/author/Wei-Guo/2312894514}{Semantic Scholar} \\
    \href{https://github.com/AlexandreGUO2001/}{Github} \\
    \href{https://www.linkedin.com/in/wei-guo-264b97293/}{LinkedIn}
}
\address{Email: \href{mailto:wei.guo@gatech.edu}{\texttt{wei.guo@gatech.edu}}}


\renewenvironment{rSection}[1]{
\sectionskip
\textcolor{LinkNavy}{\textbf{\textsc{#1}}}
\sectionlineskip
\hrule
\begin{list}{}{
\setlength{\leftmargin}{1.5em}
}
\item[]
}{
\end{list}
}



\begin{document}
\begin{rSection}{Research Summary}
I study diffusion-based generative models and sampling algorithms, with a current focus on continuous/discrete diffusion LLMs.
\end{rSection}

\begin{rSection}{Education}
{\bf \href{https://www.gatech.edu/}{Georgia Institute of Technology}, \href{https://ml.gatech.edu/}{Machine Learning Center}} \hfill {\em Aug 2023 -- Present} 
\\ Ph.D. in Machine Learning\hfill
\\ M.S. in Electrical and Computer Engineering \hfill {\em May 2026}
\\ \textbf{Advisor}: Professors \href{https://yongxin.ae.gatech.edu/}{Yongxin Chen} and \href{https://mtao8.math.gatech.edu/index.html}{Molei Tao}

{\bf \href{https://www.pku.edu.cn/}{Peking University}, \href{https://www.math.pku.edu.cn/}{School of Mathematical Sciences}} \hfill {\em Sep 2019 -- Jul 2023} 
\\ B.S. in Statistics
\\ \textbf{GPA}: 3.848/4.0; \textbf{Rank}: 5/44 \hfill
\\ \textbf{Advisor}: Professor \href{https://zcrabbit.github.io/}{Cheng Zhang} 
\end{rSection}

\begin{rSection}{Research Interests}
    \begin{itemize}[leftmargin=0in]
        \item \textbf{Machine Learning}: Diffusion-based generative AI, with a focus on continuous/discrete diffusion language models and post-training (including distillation, inference-time scaling, and reinforcement learning, etc.).
        \item \textbf{Statistics and Probability}: Theoretical analysis and practical design of sampling algorithms (Markov chain Monte Carlo, non-equilibrium (e.g., denoising diffusion, stochastic localization) sampling, or learning-based neural samplers); Applied stochastic analysis with applications to optimal transport, stochastic optimal control, and statistical physics.
    \end{itemize}
\end{rSection}


\begin{rSection}{Work Experience}
{\bf \href{https://www.nvidia.com/}{NVIDIA Corporation}, \href{https://research.nvidia.com/labs/genair/}{Fundamental Generative AI Research (GenAIR) Group}} \hfill {\em May 2026 -- Aug 2026}
\\ \textbf{Manager}: \href{http://latentspace.cc/}{Arash Vahdat}; \textbf{Mentor}: \href{https://mortezamardani.github.io/morteza/}{Morteza Mardani}
\\ Research intern; working on continuous diffusion language models
\end{rSection}


\begin{rSection}{Preprints and Publications} 
    (In reversed chronological order. * denotes equal contribution or alphabetical order.)
    \itemsep -2pt
    \nocite{*}\leavevmode\printbibliography[heading=none]
\end{rSection}



\begin{rSection}{Teaching Experience} \itemsep -2pt
\begin{itemize}[leftmargin=0in]
    \item Served as a teaching assistant for the following courses at Georgia Tech:
\end{itemize}
\href{https://ae3531a.notion.site/Course-Syllabus-Fall-2024-049faa69bedb4dd2b9b353318f52d5d6}{AE 3531: Control System Analysis and Design}. 
Lecturer: \href{https://research.gatech.edu/people/lu-gan}{Lu Gan}.
\hfill {\em Fall 2024} \\
AE 8803: Optimal Transport Theory and Applications. 
Lecturer: \href{https://yongxin.ae.gatech.edu/}{Yongxin Chen}.
\hfill {\em Spring 2025} \\
AE 3530: System Dynamics and Vibration.
Lecturer: \href{https://yongxin.ae.gatech.edu/}{Yongxin Chen}.
\hfill {\em Fall 2025} \\
AE 8803: Applied Stochastic Control and Generative Modelling.
Lecturer: \href{https://yongxin.ae.gatech.edu/}{Yongxin Chen}.
\hfill {\em Spring 2026} \\
\end{rSection}


\begin{rSection}{Community Service} \itemsep -2pt
(Numbers in parentheses indicate the number of papers reviewed. Workshop reviews not included.)
\begin{itemize}[leftmargin=0in]
    \item Served as a reviewer for the following conferences:
    
    ICLR 2024 (1), ICLR 2025 (3), NeurIPS 2025 (5), ICLR 2026 (6), ICML 2026 (6). % NeurIPS 2026 (4)
    
    \item Served as a reviewer for the following journals:
    
    IEEE Transactions on Information Theory (1), Transactions on Machine Learning Research (2).
\end{itemize}
\end{rSection}




\begin{rSection}{Awards} \itemsep -2pt
    PKU merit student in the academic year of 2019-2020 \hfill {\em Oct 2020} \\
    PKU merit student pacesetter in the academic year of 2020-2021 \hfill {\em Oct 2021} \\
    PKU merit student in the academic year of 2021-2022 \hfill {\em Oct 2022} \\
    2023 Beijing Outstanding Undergraduate Graduate Award \hfill {\em Jul 2023} \\
    Best poster award at Yale FDS conference ``Recent Advances and Future Directions for Sampling'' \\
    \phantom{.} \hfill {\em Oct 2024} \\
    Gold reviewer award at ICML 2026 \hfill {\em May 2026} \\
\end{rSection}


    
\begin{rSection}{Technical Skills} \itemsep -2pt
\begin{itemize}[leftmargin=0in]
    \item \textbf{Mathematics and Statistics}: Markov chain Monte Carlo, convex optimization, stochastic analysis, optimal transport, applied partial differential equations.
    \item \textbf{Machine learning}: Diffusion/flow-based models, large language models, reinforcement learning.
    \item \textbf{ML systems}: Familiar with PyTorch distributed training (DDP/FSDP).
    \item \textbf{Programming}: Proficient in Python, including PyTorch.
    \item \textbf{Languages}: Chinese (Mandarin and Wu dialect, native), English (proficient), French (beginner).
\end{itemize}


\end{rSection}

\fancyfoot[R]{\footnotesize Last updated: \today}

\end{document}